\section{Introduction}
\label{sec:intro}

Let $\{P_t\}_{t \geq 0}$ be a family of probability measures on $\R^d$
evolving over time. The object of interest is the lagged transport geometry
encoded by
\[
    D_j \;:=\; W_2(P_{t-j}, P_t),\qquad j = 1, 2, \ldots, L.
\]
We study the case in which this curve follows the power law
\begin{equation}
    D_j = \zeta\,j^H,\qquad \zeta > 0,\; H \in (0,1),
    \label{eq:power-law}
\end{equation}
and is observed only through noisy estimates $\Dhj$ based on $n$ samples per
lag. In log scale, the observation noise depends on the same family being
inferred, so the problem is heteroscedastic and parameter-coupled.

This coupling produces one information law for the whole paper:
\[
    \mathcal I_{n,L}(H) = n \sum_{j=1}^L j^{2H},
\]
up to scale factors depending on $\zeta$ and $\sigma_0$. The same quantity
governs estimation of $H$, the residual law, the minimax separation
boundary, and adaptivity across lag regimes. Fixed $L$ and growing $L$ are
therefore two regimes of one inferential scale.

The paper develops one inferential layer for this geometry. The two-step FWLS
estimator is asymptotically oracle-equivalent, with a joint law for the
power-law parameters and an exact oracle residual law with a feasible asymptotic
counterpart. The minimax separation boundary is of order
$\bigl(n\sum_{j \leq L} j^{2H}\bigr)^{-1/2}$, with fixed-$L$ and growing-$L$
rates as two regimes of the same information law. The procedure is adaptive in
$H$, and the unknown-scale extension separates an exact decoupled split-$F$
benchmark from a same-sample bootstrap calibration path.

\paragraph{Relation to prior work.}
Direct estimation of roughness exponents for Gaussian processes is classical
\citep{istas1997quadratic,coeurjolly2001estimating}. Here the exponent is
inferred from estimated transport discrepancies, so the noise is itself
parameter-coupled. The surrounding transport-rate literature
provides the relevant first layer \citep{fournier2015rate,weed2019sharp,
goldfeld2024limit,delbarrio2023improved}, while locally stationary drift and
adaptive inference connect through \citep{dahlhaus1997fitting,
vogt2012nonparametric,tinio2025bounds,brutsche2024similarity}.
The temporal-validity horizon law of \citet{Parreira2026May} gives one
downstream use of the pair $(H,\zeta)$: those parameters govern the structural
scale of finite memory under drift. The lag-geometry object itself is the focus
here.

Section~\ref{sec:model} defines the observation model and the information
scale. Section~\ref{sec:clt} gives the feasible inference theorem.
Section~\ref{sec:test} establishes the residual law and its response
under departures. Section~\ref{sec:lower} proves the separation lower bound,
Section~\ref{sec:adaptivity} extends it to the growing-lag regime,
Section~\ref{sec:sims} provides the empirical signatures, and
Section~\ref{sec:application} gives the unknown-scale extension.
